\documentclass[12pt,a4paper]{article}

% --------------------------------------------------
% PACOTES
% --------------------------------------------------

\usepackage[utf8]{inputenc}
\usepackage[T1]{fontenc}
\usepackage[brazil]{babel}

\usepackage{geometry}
\geometry{
    a4paper,
    left=3cm,
    right=2cm,
    top=3cm,
    bottom=2cm
}

\usepackage{setspace}
\onehalfspacing

\usepackage{indentfirst}
\setlength{\parindent}{1.25cm}

\usepackage{graphicx}
\usepackage{float}
\usepackage{booktabs}
\usepackage{array}
\usepackage{longtable}

\usepackage{xcolor}

\usepackage{hyperref}
\hypersetup{
    colorlinks=true,
    linkcolor=black,
    citecolor=black,
    urlcolor=blue
}

% Código-fonte
\usepackage{listings}

\definecolor{codegray}{rgb}{0.95,0.95,0.95}

\lstset{
    backgroundcolor=\color{codegray},
    basicstyle=\ttfamily\small,
    breaklines=true,
    frame=single,
    numbers=left,
    numberstyle=\tiny,
    showstringspaces=false
}

% --------------------------------------------------
% INFORMAÇÕES DO RELATÓRIO
% --------------------------------------------------

\title{
    \textbf{Relatório de Produção de Software}\\
    \large TRON
}

\author{
    Danilo Lins\\
    Nome do Autor 2\\
    Nome do Autor 3
}

\date{\today}


% ==================================================
% DOCUMENTO
% ==================================================

\begin{document}

% --------------------------------------------------
% CAPA
% --------------------------------------------------

\begin{titlepage}

    \centering

    {\large Universidade de Brasília / CIC}\\[0.5cm]
    {\large Ciência da Computação}\\[2cm]

    {\Large\textbf{RELATÓRIO DE PRODUÇÃO DE JOGO MULTIJOGADOR}}\\[1cm]

    {\LARGE\textbf{TRON}}\\[2cm]

    \begin{flushright}
        \begin{minipage}{0.65\textwidth}
            Relatório apresentado à disciplina
            \textit{Redes de Computadores}, como requisito
            para avaliação do desenvolvimento do trabalho final.
        \end{minipage}
    \end{flushright}

    \vfill

    Danilo Lins Castro - 251019735\\
    Eduardo Sankievicz Lima - 251035102\\
    João Pedro Barros Silva Azevedo - 251004311\\
    João Vitor da Silva Lima - 251023101

    \vfill

    Brasília -- DF\\
    2026

\end{titlepage}


% --------------------------------------------------
% RESUMO
% --------------------------------------------------

\begin{abstract}

Este relatório apresenta o desenvolvimento do software
\textit{Nome do Projeto}, descrevendo seus objetivos,
requisitos, tecnologias utilizadas, decisões de projeto,
processo de implementação e resultados obtidos.

O sistema foi desenvolvido utilizando as tecnologias
\textit{Tecnologia 1}, \textit{Tecnologia 2} e
\textit{Tecnologia 3}. Durante o desenvolvimento foram
aplicados conceitos relacionados à engenharia de software,
arquitetura de sistemas, programação e testes.

\textbf{Palavras-chave:} desenvolvimento de software;
engenharia de software; programação; sistemas.

\end{abstract}


% --------------------------------------------------
% SUMÁRIO
% --------------------------------------------------

\tableofcontents

\newpage


% ==================================================
% 1. INTRODUÇÃO
% ==================================================

\section{Introdução}

O desenvolvimento de software envolve diferentes etapas,
desde a definição do problema até a implementação e
validação da solução.

Este relatório apresenta o processo de desenvolvimento do
projeto \textit{Nome do Projeto}, desenvolvido no contexto
da disciplina \textit{Nome da Disciplina}.

O projeto tem como finalidade solucionar
\textbf{descrever brevemente o problema que o software
resolve}.

Ao longo deste relatório são apresentados os objetivos do
projeto, as tecnologias utilizadas, a metodologia adotada,
a implementação do sistema, os testes realizados e os
resultados obtidos.


% ==================================================
% 2. OBJETIVOS
% ==================================================

\section{Objetivos}

\subsection{Objetivo geral}

Desenvolver um software capaz de
\textbf{descrever aqui o objetivo principal do sistema}.

\subsection{Objetivos específicos}

Os objetivos específicos do projeto são:

\begin{itemize}
    \item implementar \textbf{funcionalidade 1};
    \item implementar \textbf{funcionalidade 2};
    \item desenvolver \textbf{funcionalidade 3};
    \item aplicar conceitos de engenharia de software;
    \item realizar testes para verificar o funcionamento
    do sistema.
\end{itemize}


% ==================================================
% 3. REQUISITOS
% ==================================================

\section{Requisitos do sistema}

\subsection{Requisitos funcionais}

Os requisitos funcionais descrevem as funcionalidades que
o sistema deve oferecer.

\begin{table}[H]
\centering
\caption{Requisitos funcionais}
\label{tab:requisitos-funcionais}

\begin{tabular}{p{2cm} p{10cm}}
\toprule
\textbf{ID} & \textbf{Descrição} \\
\midrule

RF01 & O sistema deve permitir que o usuário realize
      \textbf{ação 1}. \\

RF02 & O sistema deve permitir que o usuário realize
      \textbf{ação 2}. \\

RF03 & O sistema deve apresentar
      \textbf{informação 3}. \\

\bottomrule
\end{tabular}

\end{table}


\subsection{Requisitos não funcionais}

Entre os requisitos não funcionais definidos para o sistema
estão:

\begin{itemize}
    \item facilidade de utilização;
    \item desempenho adequado;
    \item manutenibilidade do código;
    \item segurança;
    \item compatibilidade com as plataformas definidas
    para o projeto.
\end{itemize}


% ==================================================
% 4. TECNOLOGIAS
% ==================================================

\section{Tecnologias e ferramentas}

Para o desenvolvimento do projeto foram utilizadas
diferentes tecnologias e ferramentas.

\subsection{Linguagem de programação}

O sistema foi desenvolvido utilizando
\textbf{Nome da linguagem}, escolhida devido a
\textbf{justificativa}.

\subsection{Frameworks e bibliotecas}

Foram utilizadas as seguintes bibliotecas e frameworks:

\begin{itemize}
    \item \textbf{Framework 1}: utilizado para
    \textbf{finalidade};
    
    \item \textbf{Biblioteca 2}: utilizada para
    \textbf{finalidade};
    
    \item \textbf{Biblioteca 3}: utilizada para
    \textbf{finalidade}.
\end{itemize}

\subsection{Ferramentas de desenvolvimento}

O desenvolvimento foi realizado utilizando ferramentas como
\textbf{IDE/editor}, \textbf{Git} e
\textbf{GitHub}.


% ==================================================
% 5. METODOLOGIA
% ==================================================

\section{Metodologia de desenvolvimento}

O desenvolvimento do sistema foi dividido em etapas,
permitindo organizar a implementação das funcionalidades.

Inicialmente foi realizada a análise do problema e a
definição dos requisitos. Em seguida, foi projetada a
estrutura do sistema e posteriormente realizada a
implementação.

As principais etapas foram:

\begin{enumerate}
    \item levantamento dos requisitos;
    \item definição da arquitetura;
    \item implementação;
    \item testes;
    \item correção de erros;
    \item integração das funcionalidades;
    \item validação final.
\end{enumerate}


% ==================================================
% 6. ARQUITETURA
% ==================================================

\section{Arquitetura do sistema}

A arquitetura do sistema foi definida com o objetivo de
organizar suas responsabilidades e facilitar a manutenção
do código.

\subsection{Visão geral}

O sistema é composto pelos seguintes componentes:

\begin{itemize}
    \item \textbf{Componente 1}: responsável por
    \textbf{função};
    
    \item \textbf{Componente 2}: responsável por
    \textbf{função};
    
    \item \textbf{Componente 3}: responsável por
    \textbf{função}.
\end{itemize}


\subsection{Diagrama da arquitetura}

A Figura~\ref{fig:arquitetura} apresenta uma visão geral
da arquitetura do sistema.

\begin{figure}[H]
    \centering

    % Substitua pelo arquivo da imagem
    \includegraphics[width=0.85\textwidth]{arquitetura.png}

    \caption{Arquitetura do sistema.}
    \label{fig:arquitetura}
\end{figure}


% ==================================================
% 7. IMPLEMENTAÇÃO
% ==================================================

\section{Implementação}

Nesta seção são apresentadas as principais decisões e
componentes implementados durante o desenvolvimento.

\subsection{Estrutura do projeto}

A estrutura do projeto foi organizada de forma a separar
as diferentes responsabilidades do sistema.

Por exemplo:

\begin{lstlisting}[language=Python, caption={Exemplo de estrutura do projeto}]
projeto/
|-- src/
|   |-- main.py
|   |-- models/
|   |-- services/
|   |-- controllers/
|-- tests/
|-- README.md
\end{lstlisting}


\subsection{Principais funcionalidades}

Uma das principais funcionalidades implementadas foi
\textbf{descrever funcionalidade}.

O funcionamento dessa funcionalidade ocorre da seguinte
forma:

\begin{enumerate}
    \item o usuário realiza \textbf{ação};
    \item o sistema recebe os dados;
    \item os dados são processados;
    \item o sistema realiza \textbf{operação};
    \item o resultado é apresentado ao usuário.
\end{enumerate}


\subsection{Exemplo de implementação}

A Listagem~\ref{lst:exemplo} apresenta um trecho da
implementação.

\begin{lstlisting}[language=Python,
caption={Exemplo de implementação},
label={lst:exemplo}]

def calcular(valor1, valor2):
    resultado = valor1 + valor2
    return resultado

\end{lstlisting}


% ==================================================
% 8. TESTES
% ==================================================

\section{Testes}

Foram realizados testes com o objetivo de verificar o
funcionamento das funcionalidades implementadas.

\subsection{Casos de teste}

\begin{table}[H]
\centering
\caption{Casos de teste}
\label{tab:testes}

\begin{tabular}{p{1.5cm} p{4cm} p{4cm} p{2cm}}
\toprule
\textbf{ID} &
\textbf{Entrada} &
\textbf{Resultado esperado} &
\textbf{Resultado} \\
\midrule

CT01 &
Entrada válida &
Sistema deve processar a entrada &
Passou \\

CT02 &
Entrada inválida &
Sistema deve apresentar erro &
Passou \\

CT03 &
Campo vazio &
Sistema deve rejeitar a operação &
Passou \\

\bottomrule
\end{tabular}

\end{table}


\subsection{Resultados}

Os testes realizados demonstraram que as principais
funcionalidades do sistema apresentaram o comportamento
esperado.

Caso necessário, podem ser apresentados nesta seção
gráficos, tabelas ou capturas de tela do sistema.


\begin{figure}[H]
    \centering

    % Substitua pelo arquivo da imagem
    \includegraphics[width=0.8\textwidth]{screenshot.png}

    \caption{Interface do sistema durante a execução.}
    \label{fig:screenshot}
\end{figure}


% ==================================================
% 9. DIFICULDADES E DECISÕES
% ==================================================

\section{Dificuldades e decisões de projeto}

Durante o desenvolvimento foram encontradas algumas
dificuldades relacionadas a \textbf{descrever problema}.

Para solucionar esse problema, foi adotada a estratégia
\textbf{descrever solução}.

Também foram tomadas decisões relacionadas à arquitetura,
organização do código e escolha das tecnologias.

Essas decisões contribuíram para
\textbf{descrever resultado}.


% ==================================================
% 10. RESULTADOS
% ==================================================

\section{Resultados}

Ao final do desenvolvimento, o sistema foi capaz de
\textbf{descrever resultado principal}.

As funcionalidades definidas inicialmente foram
implementadas e testadas.

A Tabela~\ref{tab:resultados} apresenta um resumo dos
resultados obtidos.

\begin{table}[H]
\centering
\caption{Resumo dos resultados}
\label{tab:resultados}

\begin{tabular}{p{7cm} p{4cm}}
\toprule
\textbf{Funcionalidade} & \textbf{Status} \\
\midrule

Funcionalidade 1 & Implementada \\
Funcionalidade 2 & Implementada \\
Funcionalidade 3 & Implementada \\
Funcionalidade 4 & Parcialmente implementada \\

\bottomrule
\end{tabular}

\end{table}


% ==================================================
% 11. CONCLUSÃO
% ==================================================

\section{Conclusão}

O desenvolvimento do projeto permitiu aplicar conceitos
relacionados à produção de software, desde a definição dos
requisitos até a implementação e realização dos testes.

O sistema desenvolvido atende aos principais objetivos
estabelecidos no início do projeto, apresentando as
funcionalidades necessárias para
\textbf{descrever finalidade}.

Além disso, o desenvolvimento possibilitou a aplicação
prática de conhecimentos relacionados a
\textbf{tecnologias/conceitos utilizados}.

Como trabalhos futuros, podem ser consideradas melhorias
como \textbf{melhoria 1}, \textbf{melhoria 2} e
\textbf{melhoria 3}.


% ==================================================
% REFERÊNCIAS
% ==================================================

\begin{thebibliography}{9}

\bibitem{exemplo}
AUTOR, Nome.

\textit{Título da obra}.

Editora, ano.

\bibitem{documentacao}
NOME DA ORGANIZAÇÃO.

\textit{Documentação oficial da tecnologia}.

Disponível em: \url{https://exemplo.com}.

Acesso em: 24 out. 2026.

\end{thebibliography}


\end{document}